\chapter{후기 및 사진}

\section{가팀}

\subsection{박준영}
언어는 결국 사람이라는 사실을, 언어학을 공부하는 우리가 가장 몰랐던 것 같습니다.

대학의 강의실, 도서관의 서가, 나의 책상에서 만나는 언어학은 그 무엇보다 편안합니다. 그 편안하지만, 영혼 없는 활자 위에서 나는 표를 그리고, 수형도를 그리고, 문제를 해결했다고 좋아합니다. 가끔 주변에 있는 사람을 모아다가 실험이라는 것도 합니다. 밀린 과제를 처리하듯 녹음을 반복하고, 계산기를 조금 두드려 보고, 귀무가설이 기각되면 좋아합니다. 학기가 끝날 즈음 보고서를 몇 장 쓰고 나면 나쁘지 않은 성적이 화면에 나타나고, 동료들은 나의 글을 칭찬해 주고, 나 스스로에게 ‘언어학도’라는 자랑스러운 명함도 쥐어주고는 합니다. 그 활자 뒤에, 계산기 뒤에 있는 사람들을 잊은 채로요.

이번 언어조사는 타자기도, 실험실도, 계산기도 아닌 ‘사람’의 언어를 다시금 찾는 시간이었습니다. 사람의 언어는 활자처럼 완벽하지도, 실험처럼 정제되어 있지도, 수학처럼 딱딱 맞아떨어지지도 않습니다. 언어조사의 과정 역시 그러했습니다. 처음 뵙는 어르신들께 어떻게든 말을 붙여야 했던 때의 아득함을 기억합니다. 성함도 모르는 어르신께 인천상륙작전 무용담을 들어야 했던 때의 혼란함을 기억합니다. 어르신께 동의서를 제대로 설명하지 못해 조사가 엎어질 뻔했을 때의 초조함을 기억합니다. 이틀 간의 현장실습은 예측할 수 없는 사건의 연속이었습니다.

그럼에도, 제 마음속에는 따뜻함이 더 많이 묻어 있는 듯합니다. 기쁨과 열의에 차 적극적으로 설명해 주시던 어르신, 사례비를 한사코 거절하시던 어르신, 학생들 공부에 도움 되는 것이라면 뭐든 좋다고 말씀하시던 어르신, 언어조사 끝나면 꼬막 축제에 갈 것이라 웃으며 말씀하시던 어르신…. 이 모든 기억이 언어조사를 계속하는 동력이 되는 것 같습니다.

말소리에 배어 있는 인간미도 언어조사의 큰 매력이라고 하겠습니다. 하루의 언어조사를 마친 뒤 회의를 열어 각 팀의 조사 결과를 공유했는데, 화자마다 방언형 차이는 물론, 의미 영역이나 음운 체계에도 미묘한 편차가 있다는 사실을 알게 되었습니다. 이 편차에 대해 토론하고 언어를 어떻게 기술할 것인지에 대한 합의점을 찾아나가는 과정이 과연 ‘현장언어학’답다고 느껴 대단히 매력적이었습니다.

“잊혀 가는 것을 위해 공부하고 싶습니다.” 내가 고등학생 시절 언어학과로의 진학을 꿈꾸며 자기소개서에 무심코 담았던 문장입니다. 이 문장을 쓰고서 사 년이 지난 뒤에야, ‘잊혀 가는 것’에 다가가는 첫 걸음을 뗀 것 같아 조금은 겸연쩍습니다. 뱉은 말을 지키려면 오랜 시간이 걸리겠지요. 흡족할 때까지 잊혀 가는 것, 잊혀 가는 사람들에 대해 생각해 보아야겠습니다. 그제서야 언어는 결국 사람이라는 사실을, 완전히 이해할지도 모르겠습니다.

\subsection{유승민}
언어학도로서 안락의자에서 벗어나 진짜 언어학을 배운 뜻깊은 시간이었다. 입학 이후에 언어학이 왜 인문학인지에 의문을 가진 경험이 더러 있었다. 하지만 이번 기회에 직접 자료제공인을 만나며, 어렴풋이나마 우리가 공부하는 것은 언어가 아니라 언어를 사용하는 사람임을, 트리를 그리는 것도, 표를 그리는 것도 이를 위한 부수적인 수단임을 새삼 깨달았다. 자료조사를 준비하고, 참여하는 과정 또한 즐거웠다.

첫 수업부터 흔히 우리가 수업에서 접하는 것과는 사뭇 다른 내용을 배우겠다는 것은 알았다. 물론 당연히 선생님이 계시고 강의자료가 있었지만, 그 내용은 흔히 접하는 사변적인 이론이나 실험의 결과, 개념적 사실보다는 마치 워크숍애서 실무를 배우는 것과 같았다. 그리고, 몇 주의 수업 이후에 나는 수업에서 배운 조금의 잔상 이후에 벌교로 향하는 기차에 올라탔다. 트리를 그릴 줄만 알던 내가 언어조사에 임하는 상황이 익숙지 않았다. 첫날에 직접 자료제공인을 구하고, 적합한 공간을 물색하고, 자료제공인과 동료들과 함꼐 언어조사를 하는 과정은 실제로도 참 힘들었다. 그렇지만 자료조사를 마친 후에 내일의 자료조사를 논의하는 숙소에서 새삼 하루 사이에 많은 경험을 했다고 느꼈다. 실제 언어는 책 밖에 존재하고, 언어를 배우는 것은 언어조사가 시작이자 끝이다. 교실이나 자취방에서 벗어나 밖에서 진짜 언어학을 할 기회가 주어져서 참 다행이라고 생각한다. 언어조사의 경험이 향후의 학업 수행에 큰 의의를 가질 것임을 믿어 의심치 않는다.

물론 힘들겠지만, 그리고 보고서의 작성이 끝나지 않은 현재 시점에서 실제로 아직도 힘들지만, 아직 언어조사의 경험이 없는 학우들에게 이 수업을 추천하는 이유이다. 수강생이 몰리면 지금보다는 덜 재밌어지겠지만, 본 과목이 전공필수 과목을 지정되기를 바란다.

\subsection{이채현}
졸업 전에 언어학도로서 현장에서 언어를 경험하고 기록하는 작업을 해보고싶어서 수업을 신청하였고, 벌교로 떠났던 언어조사는 제 기대만큼이나 무척 즐겁고 유익한 시간이었습니다. 소규모 인원으로 조사를 준비하고 진행하였는데, 소규모다 보니까 팀 전체가 함께 소통할 수 있어서 좋았습니다. 또한 버스가 아닌, 기차를 이용하고 렌터카를 빌려서 벌교를 다녀올 수 있었던 점도 특별했습니다. 운전을 비롯하여 여러 도움을 주신 강민하 조교님께 감사드립니다. 사실 언어조사를 준비하는 과정은 생각보다 맨 땅에 헤딩하는 느낌이었습니다. 작년에 했던 형식이 참고용으로 존재하긴 했으나, 주어진 가이드라인이 명확하게 있다기보다는 현장에 갑자기 뛰어들게 된 느낌이 없지는 않았던 거 같습니다. 서울 태생에 서울 화자인 저로서는 더더욱 방언의 특성을 공부하고 소리로 언어적 특성을 구분하는 것이 생각보다 어려웠습니다. 특히 음성학에서 했던 음성 기호를 까먹어서 다시 IPA 차트를 찾아보고 있는 저를 발견할 수 있었습니다. 아직 정밀 전사 작업과 분석이 남기는 했으나, 어찌됐든 언어조사 과정은 제게 무척이나 즐겁고 특별한 기억으로 남았습니다. 졸업 전에 이 수업을 안 들었으면 후회했을 거같습니다. 벌교 지역의 어르신들과 이야기 나누는 것도 즐거웠고, 열심히 준비해 간 질문지를 하나씩 클리어해가는 것에서 묘한 쾌감을 느끼기도 했습니다. 제보자분들과의 대화 속에서 저희가 사전조사로 형태를 예측했던 어휘가 그대로 나올 때, 셋다 흥분해서 종이를 막 넘기고 펜을 끄적이곤 했는데, 그게 되게 기억에 남는 거같아요. 화자를 모시고 조사 환경을 세팅하는 데 조금 우여곡절을 겪기도 했으나, 돌이켜 보면 그 또한 유쾌한 에피소드인 거같습니다. 저희를 향한 어르신들의 관심이 이렇게 크다니, 사실 그에 감사한 마음을 가져야 하는 것도 같아요. 손주 손녀들처럼 대해주시고 반갑게 맞아주신 벌교 지역 어르신들께 감사하고, 또 이 모든 여정을 통솔해주신 교수님께도 감사하다는 말씀 드리고 싶습니다.

다만 아쉬운 점은, 처음부터 반장을 정해서 전체 일정을 정하고 조금 더 체계적으로 준비했으면 좋았을 거같습니다. 또한 반장이 해야할 일에 대한 가이드라인이 명확했으면 좋겠습니다. 조교님께서 감사하게도 많은 부분을 도와주셨는데, 이 부분이 인수인계 형태로 반장이 해야 할일로 정리되어 전달된다면 학생 주도의 조사 준비가 더 원활할 거같습니다. 지금은 뭘 어떻게 준비해야 되는지 정확히 모르는 상황에서 학생들끼리 스스로 해야 하는 상황이라 방향성을 잡기 어려운 거같습니다.

\section{나팀}

\subsection{이상인}
저는 고향도 서울이고 국내여행도 좀처럼 가지 않는 터라, 지방에 내려가서 이렇게 사투리를 많이 들을 일이 좀처럼 없었습니다. 그래서 사실 답사 당일까지도, 너무 긴장되고 제 자신이 너무 작게 느껴졌습니다. 방언 공동체를 어떻게 다뤄야 할지도 전혀 몰랐고요. 그럼에도 답사는 아주 즐거운 경험이었습니다. 어르신들께서는 생각보다 저희를 반겨 주셨고, 복지관이나 숙소 등 시설은 부족한 점이 전혀 없었으며, 조사 자체도 대단히 즐겁고 수월하게 진행되었습니다. 또, 돈을 모아서 어르신들께 선물도 드리고 밤에 간식도 먹고 하니 우리 스스로의 돈으로 만들어간다는 좋은 기분이 들면서도, 학과에서 교통(KTX), 숙소, 식사까지 지원을 해 주니 금전적으로 큰 부담이 되지 않았습니다. 꼬막도 너무 맛있었고요!! 특히 이번 언어조사는 지금까지와는 달리 매우 적은 인원(6명..)으로 진행되어, 진행상의 편의도 있었고 수강생끼리도 더 끈끈해지지 않았나 하는 생각이 듭니다. 물론 이제 압도적인 양의 전사와, 검토와, 2차 검토와, 보고서 작성이 남아있지만... 이 수업을 한 번 더 듣고 싶냐 하면 저는 그렇다고 대답할 것 같습니다. 노동을 (좀 많이..) 수반하기는 하지만 언어학과가 보내주는 엄청 재밌는 컨텐츠가 포함된 여행이니까요!

\subsection{문준모}
학기가 시작하고 수업 첫 날, 저 이외에 네 명밖에 없는 강의실을 보고는 솔직히 많이 당황했습니다. 이후 두 명이 늘어 수강생 여섯과 조교님 둘, 그리고 최계영 선생님과 함께 떠나는 언어 조사가 되었습니다만, KTX를 타고 현지에 도착해, 9인승 카니발에서 밴드 음악을 들으며 창문 밖 남도의 정취에 젖어가던 와중에도 제 마음은 설렘 반, 당혹감 반이었던 것 같습니다.

마침내 조사 장소에 들어섰을 때, 저희는 이상적인 제보자를 찾고, 적절한 공간을 물색하며 장비 세팅을 마치는 데 각기 바삐 움직였습니다. 솔직히 말하면 첫 날 조사 당시의 기억은 흐릿하게밖에 남아 있지 않습니다. 선명하게 기억나는 것은 어르신들께서 떠나신 뒤 팀원끼리 정자에 모여 앉아, 질문지에 휘갈겨 쓴 글씨를 두고 이야기꽃을 피운 일입니다. 그제서야 실감했습니다. 아, 이렇게 내가 하루의 언어 조사를 마쳤구나!

그 이후의 일정은 줄곧 순조로웠습니다. 벌교의 명물이라는 꼬막 정식을 먹으면서도, 숙소에 모여 두 시간여의 열띤 회의를 하면서도, 다음날 다시 김정숙(화자7) 어르신을 맞이했을 때도 즐거운 추억만 가득 쌓았습니다. 질문지의 빈칸을 채워야 한다는 압박감은 사라지지 않았지만, 적어도 그 과정이 어디까지나 사람을 대하고, 이야기에 귀기울이고, 울고 웃는 것이라는, 수업에서도 배웠던 사실을 진정으로 깨달을 수 있었습니다.

여전히 저는 미숙한 조사자입니다. 앞으로 한 번, 두 번, 세 번의 조사를 더 떠난다고 해서 노련미를 뽐낼 수 있으리라 생각되지는 않습니다. 또 전사와 보고서 취합이라는, 현장 조사만큼이나 중요한 과업에 있어는 시작선에 겨우 섰을 따름입니다. 그러나 만일 본 수업을 다시 듣고 싶으냐고 묻는다면 저는 그렇다고 대답하겠습니다. 이제까지의 긍정적 경험을 토대로 생각합니다. 아, 이것이 ≪언어 조사와 분석≫이구나!

\subsection{전지민}
아직 언어학을 깊게 배우지 않은 타과생의 입장에서 이 수업이 제가 언어학에 보다 흥미를 가질 수 있게 해주는 계기가 되었다고 생각합니다. 단순히 재미있어 보인다는 이유만으로 언어학 다전공을 신청하고, 또 단순히 재미있어 보인다는 이유만으로 이 수업을 신청하였지만 언어조사 준비부터 실제 현장에서의 언어조사까지 어느 하나도 새롭고 흥미로운 경험이었습니다. 아마 제가 수강생 중 이 과목을 수강하기 위해 필요한 사전 지식이 가장 적었을 것 같은데, 부족한 저를 도와주시고 이 언어조사를 성공적으로 이끌어 주신 학우분들 정말 감사드립니다. 

언어조사를 진행하며 느낀 점은 운이 좋게도 저희는 제보자 분들이 굉장히 협조적이었던 것 같습니다. 저희 조의 경우 1일차에 조사한 제보자 중 한 분이 굉장히 열정적으로 답변을 해 주시고 다음 날까지 참석해 주신 것이 인상 깊었고, 다른 조의 경우에도 제보자 분들이 굉장히 열정적으로 자신들의 이야기를 해 주셨다고 들었습니다. 덕분에 벌교의 말에 대해 많이 알아갈 수 있었습니다. 교수님께서도 강조하셨지만 언어 조사라는 것은 단순한 학술적 탐구가 아닌 사람을 상대하는 일인데, 이런 식으로 사람을 상대하며 배우는 과정이 언어학과에서도 희귀한 경험이겠지만, 자연과학대학 소속인 저에게는 정말로 이색적이고 특별한 경험으로 남을 것 같습니다. 매일같이 실험실에서 보내던 시간 대신, 학우들, 조교님, 교수님, 그리고 제보자분들과 함께 말을 배우고, 또 말을 배우는 과정을 배우는 경험을 할 수 있는 기회를 받아 정말 감사했습니다. 1학기 때 음성학 수업을 들으며 P-side는 저와 맞지 않는다는 생각을 잠깐 했는데 이 수업을 계기로 조금 더 공부하고 싶어졌습니다. 남은 학부 기간 동안 최대한 언어학에 도전해보겠습니다. 